\section{Very ample sheaves}
\label{section-very-ample}

\noindent
Recall that given a quasi-coherent sheaf $\mathcal{E}$ on a scheme
$S$ the {\it projective bundle} associated to $\mathcal{E}$ is the morphism
$\mathbf{P}(\mathcal{E}) \to S$, where
$\mathbf{P}(\mathcal{E}) = \underline{\text{Proj}}_S(\text{Sym}(\mathcal{E}))$,
see
Constructions, Definition \ref{constructions-definition-projective-bundle}.

\begin{definition}
\label{definition-very-ample}
Let $f : X \to S$ be a morphism of schemes.
Let $\mathcal{L}$ be an invertible $\mathcal{O}_X$-module.
We say $\mathcal{L}$ is {\it relatively very ample} or more
precisely {\it $f$-relatively very ample}, or
{\it very ample on $X/S$}, or {\it $f$-very ample} if
there exist a quasi-coherent $\mathcal{O}_S$-module
$\mathcal{E}$ and an immersion $i : X \to \mathbf{P}(\mathcal{E})$
over $S$ such that
$\mathcal{L} \cong i^*\mathcal{O}_{\mathbf{P}(\mathcal{E})}(1)$.
\end{definition}



\noindent
Since there is no assumption of quasi-compactness in this definition it is not
true in general that a relatively very ample invertible sheaf is a relatively
ample invertible sheaf.

\begin{lemma}
\label{lemma-ample-very-ample}
\begin{reference}
\cite[II, Proposition 4.6.2]{EGA}
\end{reference}
Let $f : X \to S$ be a morphism of schemes.
Let $\mathcal{L}$ be an invertible $\mathcal{O}_X$-module.
If $f$ is quasi-compact and $\mathcal{L}$ is a relatively
very ample invertible sheaf, then $\mathcal{L}$ is a relatively
ample invertible sheaf.
\end{lemma}

\begin{proof}
By definition there exists quasi-coherent $\mathcal{O}_S$-module
$\mathcal{E}$ and an immersion $i : X \to \mathbf{P}(\mathcal{E})$
over $S$ such that
$\mathcal{L} \cong i^*\mathcal{O}_{\mathbf{P}(\mathcal{E})}(1)$.
Set $\mathcal{A} = \text{Sym}(\mathcal{E})$, so
$\mathbf{P}(\mathcal{E}) = \underline{\text{Proj}}_S(\mathcal{A})$
by definition. The graded $\mathcal{O}_S$-algebra $\mathcal{A}$
comes equipped with a map
$$
\psi :
\mathcal{A} \to
\bigoplus\nolimits_{n \geq 0}
\pi_*\mathcal{O}_{\mathbf{P}(\mathcal{E})}(n) \to
\bigoplus\nolimits_{n \geq 0}
f_*\mathcal{L}^{\otimes n}
$$
where the second arrow uses the identification
$\mathcal{L} \cong i^*\mathcal{O}_{\mathbf{P}(\mathcal{E})}(1)$.
By adjointness of $f_*$ and $f^*$ we get a morphism
$\psi : f^*\mathcal{A} \to \bigoplus_{n \geq 0}\mathcal{L}^{\otimes n}$.
We omit the verification that the morphism $r_{\mathcal{L}, \psi}$
associated to this map is exactly the immersion $i$.
Hence the result follows from
part (6) of Lemma \ref{lemma-characterize-relatively-ample}.
\end{proof}

\noindent
To arrive at the correct converse of this lemma we ask
whether given a relatively ample
invertible sheaf $\mathcal{L}$ there exists an integer $n \geq 1$ such
that $\mathcal{L}^{\otimes n}$ is relatively very ample? In general this
is false. There are several things that prevent this from being true:
\begin{enumerate}
\item Even if $S$ is affine, it can happen that no finite integer
$n$ works because $X \to S$ is not of finite type, see
Example \ref{example-not-finite-type-proj}.
\item The base not being quasi-compact means the result can be
prevented from being true even with $f$ finite type. Namely, given
a field $k$ there exists a scheme $X_d$ of finite type over $k$ with
an ample invertible sheaf $\mathcal{O}_{X_d}(1)$ so that the smallest
tensor power of $\mathcal{O}_{X_d}(1)$ which is very ample is the $d$th
power. See Example \ref{example-not-bounded}.
Taking $f$ to be the disjoint union of the schemes $X_d$ mapping
to the disjoint union of copies of $\Spec(k)$ gives an example.
\end{enumerate}
To see our version of the converse take a look at
Lemma \ref{lemma-finite-type-ample-very-ample} below.
We will do some preliminary work before proving it.

\begin{example}
\label{example-very-ample}
Let $S$ be a scheme. Let $\mathcal{A}$ be a quasi-coherent graded
$\mathcal{O}_S$-algebra generated by $\mathcal{A}_1$ over $\mathcal{A}_0$.
Set $X = \underline{\text{Proj}}_S(\mathcal{A})$. In this case
$\mathcal{O}_X(1)$ is invertible and very ample on $X/S$. Namely,
the morphism associated to the graded $\mathcal{O}_S$-algebra map
$$
\text{Sym}_{\mathcal{O}_X}^*(\mathcal{A}_1)
\longrightarrow
\mathcal{A}
$$
is a closed immersion $X \to \mathbf{P}(\mathcal{A}_1)$ which pulls back
$\mathcal{O}_{\mathbf{P}(\mathcal{A}_1)}(1)$ to $\mathcal{O}_X(1)$, see
Constructions, Lemma
\ref{constructions-lemma-surjective-generated-degree-1-map-relative-proj}.
\end{example}

\begin{example}
\label{example-not-finite-type-proj}
Let $k$ be a field.
Consider the graded $k$-algebra
$$
A = k[U, V, Z_1, Z_2, Z_3, \ldots]/I
\quad
\text{with}
\quad
I = (U^2 - Z_1^2, U^4 - Z_2^2, U^6 - Z_3^2, \ldots)
$$
with grading given by $\deg(U) = \deg(V) = \deg(Z_1) = 1$
and $\deg(Z_d) = d$.
Note that $X = \text{Proj}(A)$ is covered by $D_{+}(U)$ and
$D_{+}(V)$. Hence the sheaves $\mathcal{O}_X(n)$ are all
invertible and isomorphic to $\mathcal{O}_X(1)^{\otimes n}$.
In particular $\mathcal{O}_X(1)$ is ample and $f$-ample
for the morphism $f : X \to \Spec(k)$.
We claim that no power of $\mathcal{O}_X(1)$ is $f$-relatively very ample.
Namely, it is easy to see that $\Gamma(X, \mathcal{O}_X(n))$
is the degree $n$ summand of the algebra $A$. Hence if $\mathcal{O}_X(n)$
were very ample, then $X$ would be a closed subscheme of a projective
space over $k$ and hence of finite type over $k$. On the other hand
$D_{+}(V)$ is the spectrum of
$k[t, t_1, t_2, \ldots]/(t^2 - t_1^2, t^4 - t_2^2, t^6 - t_3^2, \ldots)$
which is not of finite type over $k$.
\end{example}

\begin{example}
\label{example-not-bounded}
Let $k$ be an infinite field. Let $\lambda_1, \lambda_2, \lambda_3, \ldots$
be pairwise distinct elements of $k^*$. (This is not strictly necessary,
and in fact the example works perfectly well even if all $\lambda_i$
are equal to $1$.)
Consider the graded $k$-algebra
$$
A_d = k[U, V, Z]/I_d
\quad
\text{with}
\quad
I_d = (Z^2 - \prod\nolimits_{i = 1}^{2d} (U - \lambda_i V)).
$$
with grading given by $\deg(U) = \deg(V) = 1$ and $\deg(Z) = d$.
Then $X_d = \text{Proj}(A_d)$ has ample invertible sheaf
$\mathcal{O}_{X_d}(1)$. We claim that if $\mathcal{O}_{X_d}(n)$
is very ample, then $n \geq d$. The reason for this is that $Z$
has degree $d$, and hence $\Gamma(X_d, \mathcal{O}_{X_d}(n)) =
k[U, V]_n$ for $n < d$. Details omitted.
\end{example}

\begin{lemma}
\label{lemma-relatively-very-ample-separated}
Let $f : X \to S$ be a morphism of schemes.
Let $\mathcal{L}$ be an invertible sheaf on $X$.
If $\mathcal{L}$ is relatively very ample on $X/S$ then
$f$ is separated.
\end{lemma}

\begin{proof}
Being separated is local on the base (see
Schemes, Section \ref{schemes-section-separation-axioms}).
An immersion is separated
(see Schemes, Lemma \ref{schemes-lemma-immersions-monomorphisms}).
Hence the lemma follows since locally $X$ has an immersion into
the homogeneous spectrum of a graded ring which is separated, see
Constructions, Lemma \ref{constructions-lemma-proj-separated}.
\end{proof}

\begin{lemma}
\label{lemma-relatively-very-ample}
Let $f : X \to S$ be a morphism of schemes.
Let $\mathcal{L}$ be an invertible sheaf on $X$.
Assume $f$ is quasi-compact. The following are
equivalent
\begin{enumerate}
\item $\mathcal{L}$ is relatively very ample on $X/S$,
\item there exists an open covering $S = \bigcup V_j$ such
that $\mathcal{L}|_{f^{-1}(V_j)}$ is relatively very ample
on $f^{-1}(V_j)/V_j$ for all $j$,
\item there exists a quasi-coherent sheaf of graded
$\mathcal{O}_S$-algebras $\mathcal{A}$ generated in degree
$1$ over $\mathcal{O}_S$ and a map of graded $\mathcal{O}_X$-algebras
$\psi : f^*\mathcal{A} \to \bigoplus_{n \geq 0} \mathcal{L}^{\otimes n}$
such that $f^*\mathcal{A}_1 \to \mathcal{L}$ is surjective and the
associated morphism
$r_{\mathcal{L}, \psi} : X \to \underline{\text{Proj}}_S(\mathcal{A})$
is an immersion, and
\item $f$ is quasi-separated, the canonical map
$\psi : f^*f_*\mathcal{L} \to \mathcal{L}$ is surjective, and
the associated map $r_{\mathcal{L}, \psi} : X \to \mathbf{P}(f_*\mathcal{L})$
is an immersion.
\end{enumerate}
\end{lemma}

\begin{proof}
It is clear that (1) implies (2). It is also clear that
(4) implies (1); the hypothesis of quasi-separation
in (4) is used to guarantee that $f_*\mathcal{L}$ is quasi-coherent via
Schemes, Lemma \ref{schemes-lemma-push-forward-quasi-coherent}.

\medskip\noindent
Assume (2). We will prove (4).
Let $S = \bigcup V_j$ be an open covering as in (2).
Set $X_j = f^{-1}(V_j)$ and $f_j : X_j \to V_j$ the
restriction of $f$. We see that $f$ is separated by
Lemma \ref{lemma-relatively-very-ample-separated} (as being
separated is local on the base). By assumption there exists a
quasi-coherent $\mathcal{O}_{V_j}$-module $\mathcal{E}_j$ and an immersion
$i_j : X_j \to \mathbf{P}(\mathcal{E}_j)$ with
$\mathcal{L}|_{X_j} \cong i_j^*\mathcal{O}_{\mathbf{P}(\mathcal{E}_j)}(1)$.
The morphism $i_j$ corresponds to a surjection
$f_j^*\mathcal{E}_j \to \mathcal{L}|_{X_j}$, see
Constructions, Section \ref{constructions-section-projective-bundle}.
This map is adjoint to a map $\mathcal{E}_j \to f_*\mathcal{L}|_{V_j}$
such that the composition
$$
f_j^*\mathcal{E}_j \to (f^*f_*\mathcal{L})|_{X_j} \to \mathcal{L}|_{X_j}
$$
is surjective. We conclude that
$\psi : f^*f_*\mathcal{L} \to \mathcal{L}$ is surjective. Let
$r_{\mathcal{L}, \psi} : X \to \mathbf{P}(f_*\mathcal{L})$ be the
associated morphism. We still have to show that $r_{\mathcal{L}, \psi}$
is an immersion; we urge the reader to prove this for themselves.
The $\mathcal{O}_{V_j}$-module map $\mathcal{E}_j \to f_*\mathcal{L}|_{V_j}$
determines a homomorphism on symmetric algebras, which in turn defines a
morphism
$$
\mathbf{P}(f_*\mathcal{L}|_{V_j}) \supset U_j \longrightarrow
\mathbf{P}(\mathcal{E}_j)
$$
where $U_j$ is the open subscheme of
Constructions, Lemma \ref{constructions-lemma-morphism-relative-proj}.
The compatibility of $\psi$ with $\mathcal{E}_j \to f_*\mathcal{L}|_{V_j}$
shows that $r_{\mathcal{L}, \psi}(X_j) \subset U_j$ and that there is a
factorization
$$
\xymatrix{
X_j \ar[r]^-{r_{\mathcal{L}, \psi}} & U_j \ar[r] &
\mathbf{P}(\mathcal{E}_j)
}
$$
We omit the verification.
This shows that $r_{\mathcal{L}, \psi}$ is an immersion.

\medskip\noindent
At this point we see that (1), (2) and (4) are equivalent.
Clearly (4) implies (3). Assume (3). We will prove (1).
Let $\mathcal{A}$ be a quasi-coherent sheaf of graded $\mathcal{O}_S$-algebras
generated in degree $1$ over $\mathcal{O}_S$. Consider the map of
graded $\mathcal{O}_S$-algebras $\text{Sym}(\mathcal{A}_1) \to \mathcal{A}$.
This is surjective by hypothesis and hence induces a closed immersion
$$
\underline{\text{Proj}}_S(\mathcal{A})
\longrightarrow
\mathbf{P}(\mathcal{A}_1)
$$
which pulls back $\mathcal{O}(1)$ to $\mathcal{O}(1)$,
see Constructions, Lemma
\ref{constructions-lemma-surjective-generated-degree-1-map-relative-proj}.
Hence it is clear that (3) implies (1).
\end{proof}

\begin{lemma}
\label{lemma-very-ample-base-change}
Let $f : X \to S$ be a morphism of schemes.
Let $\mathcal{L}$ be an invertible $\mathcal{O}_X$-module.
Let $S' \to S$ be a morphism of schemes.
Let $f' : X' \to S'$ be the base change of $f$ and denote
$\mathcal{L}'$ the pullback of $\mathcal{L}$ to $X'$.
If $\mathcal{L}$ is $f$-very ample, then $\mathcal{L}'$ is $f'$-very ample.
\end{lemma}

\begin{proof}
By Definition \ref{definition-very-ample} there exists there exist a
quasi-coherent $\mathcal{O}_S$-module $\mathcal{E}$ and an immersion
$i : X \to \mathbf{P}(\mathcal{E})$ over $S$ such that
$\mathcal{L} \cong i^*\mathcal{O}_{\mathbf{P}(\mathcal{E})}(1)$.
The base change of $\mathbf{P}(\mathcal{E})$ to $S'$ is
the projective bundle associated to the pullback $\mathcal{E}'$
of $\mathcal{E}$ and the pullback of
$\mathcal{O}_{\mathbf{P}(\mathcal{E})}(1)$
is
$\mathcal{O}_{\mathbf{P}(\mathcal{E}')}(1)$, see
Constructions, Lemma \ref{constructions-lemma-relative-proj-base-change}.
Finally, the base change
of an immersion is an immersion
(Schemes, Lemma \ref{schemes-lemma-base-change-immersion}).
\end{proof}







